\begin{textsample}{POS Dim 1 – se – Score 38.00 – se/se\_em\_16.txt}  \label{ex:f1_pos_114}
3.4 Ciências Humanas e Sociais Aplicadas As aprendizagens essenciais definidas na BNCC do Ensino Médio estão organizadas por áreas do conhecimento ( Linguagens e suas Tecnologias, Matemática e suas Tecnologias, Ciências da Natureza e suas Tecnologias, Ciências Humanas e Sociais Aplicadas ), conforme estabelecido no artigo 35 -A da LDB.
Desde \textbf{que} foram introduzidas nas DCNEM/1998 ( Parecer CNE/CEB nº 15/1998 ), as áreas do conhecimento têm por finalidade integrar dois ou mais componentes do currículo, para melhor compreender a complexa realidade e atuar nela.
Essa concepção, fortalecida no Parecer CNE/CP nº 11/2009, \textbf{que} trata de inovação curricular : > “ (... ) não exclui necessariamente as \textbf{disciplinas}, com suas especificidades e saberes próprios \textbf{historicamente} construídos, mas, sim, \textbf{implica} o fortalecimento das relações entre elas e a sua \textbf{contextualização} para apreensão e intervenção na realidade, requerendo trabalho conjugado e cooperativo dos seus professores no planejamento e na execução dos planos de ensino ” ( Parecer CNE/CP nº 11/2009, p.08 ).
As Ciências Humanas, tanto no Ensino Fundamental como no Ensino Médio, apresentam aprendizagens \textbf{focadas} no desenvolvimento das competências de identificar, analisar, comparar e interpretar ideias, pensamentos, fenômenos e processos históricos, geográficos, sociais, econômicos, políticos e culturais.
Essas competências conduzem os estudantes a elaborar hipóteses, construir argumentos e atuar no mundo, lançando mão dos objetos de conhecimento dos componentes da área.
No Ensino Médio, com a incorporação da Filosofia e da Sociologia, a área de Ciências Humanas e Sociais Aplicadas ( CHSA ) propõe o aprofundamento e a ampliação da base \textbf{conceitual} e dos modos de construção da argumentação e \textbf{sistematização} do \textbf{raciocínio}, operacionalizados com base em procedimentos analíticos e interpretativos.
Nessa etapa, como a aprendizagem está centrada nas experiências dos jovens \textbf{enquanto} protagonistas, deve -se estimular \textbf{uma} leitura de mundo pautada numa visão crítica e contextualizada da realidade, no domínio \textbf{conceitual} e na elaboração e aplicação de interpretações sobre as relações, os processos e as múltiplas dimensões da existência humana.
Dessa forma, a BNCC da área de CHSA está organizada de modo a tematizar e problematizar algumas \textbf{categorias} da área, fundamentais à formação dos estudantes : Tempo e Espaço ; Territórios e Fronteiras ; Indivíduo, Natureza, Sociedade, Cultura e Ética ; e Política e Trabalho.
Essa empreitada pode ser operacionalizada a partir da \textbf{problematização} numa perspectiva \textbf{teórica} plural \textbf{que} afaste as interpretações dogmáticas e unilaterais da realidade social, e \textbf{que} contribua para \textbf{um} melhor entendimento da diversidade sociocultural.
O entendimento da referida área nesta dimensão pluralista relaciona -se com as orientações e diretrizes \textbf{que} nortearam a construção dos documentos \textbf{que} estão propondo \textbf{uma} reorganização do Ensino Médio, destacando -se o \textbf{compromisso} com a formação da cidadania, \textbf{que} \textbf{implica} a compreensão da realidade econômica, social e política na qual o indivíduo está inserido, oferecendo, a partir da análise, condições \textbf{que} propiciem \textbf{uma} atuação transformadora.
Em CHSA, quando colocamos em evidência a \textbf{problematização}, estamos nos referindo a \textbf{uma} prática pedagógica pautada na \textbf{contextualização} ; ambas ( \textbf{problematização} e \textbf{contextualização} ) se complementam na construção do conhecimento filosófico, geográfico, histórico e sociológico ; são saberes situados a partir do conhecimento elaborado com \textbf{um} dado recorte, sempre segundo alguém e em determinado contexto.
Assim, num sentido mais amplo, politiza -se o conhecimento e permite -se \textbf{que} os estudantes promovam \textbf{um} diálogo entre si e com os conhecimentos \textbf{historicamente} construídos.
Como \textbf{uma} das ferramentas para mobilizar esses saberes nas dimensões cognitiva e socioemocional, as metodologias ativas, mescladas com \textbf{uma} parte expositiva, se mostram como aliadas \textbf{enquanto} práticas pedagógicas estruturadas, de modo \textbf{que} os estudantes participem do seu processo de aprendizado, sejam estimulados a enfrentarem situações complexas e resolverem problemas práticos, contribuindo para o desenvolvimento da criatividade, do pensamento crítico por meio da autonomia, argumentação, responsabilidade, proatividade, colaboração, diálogo respeitoso, trabalho em equipe e independência.
De acordo com o Caderno de Práticas da BNCC, trabalhar com metodologias ativas – colaborativas e cooperativas, permite a construção do conhecimento através do desenvolvimento de importantes competências, como : - Saber buscar e investigar informações com criticidade ( critérios de seleção e priorização ) a fim de atingir determinado objetivo, a partir da formulação de perguntas ou de desafios dados pelos educadores ; - Compreender a informação, analisando -a em diferentes níveis de complexidade, contextualizando -a e associando -a a outros conhecimentos ; - Interagir, negociar e comunicar -se com o grupo, em diferentes contextos e momentos ; - Conviver e agir com inteligência emocional, identificando e desenvolvendo atitudes positivas para a aprendizagem colaborativa ; - Ter autogestão afetiva, reconhecendo atitudes interpessoais facilitadoras e dificultadoras para a qualidade da aprendizagem, lidando com o \textbf{erro} e as frustrações, e sendo flexível ; - Tomar decisão individualmente e em grupo, avaliando os pontos positivos e negativos envolvidos ; - Desenvolver a capacidade de liderança ; - Resolver problemas, executando \textbf{um} projeto ou \textbf{uma} ação e propondo soluções. \textbf{Dito} isso, acreditamos \textbf{que}, no \textbf{atual} contexto da sociedade \textbf{globalizada}, é também no tempo e no espaço da Educação Básica \textbf{que} valores como cidadania, consciência ecológica, direitos humanos, democracia e solidariedade, por exemplo, devem ser analisados e vivenciados pelos estudantes.
São princípios \textbf{que} permitem desnaturalizar e romper com os círculos de desigualdade e de preconceitos \textbf{que} ainda dividem e agridem a \textbf{humanidade} e, em particular, a sociedade brasileira.
De \textbf{um} modo mais abrangente, o Novo Ensino Médio traz a oportunidade do estudante refletir sobre si mesmo, seus desejos, anseios, potencialidades, por meio de \textbf{um} trabalho pedagógico intencional e sistemático a partir de conhecimentos e experiências \textbf{que} desenvolvem a determinação, esforço, autoeficácia, perseverança, autoavaliação, compreensão sobre o mundo do trabalho, isto é, as dimensões pessoal, cidadã e profissional.
Nesse sentido, os jovens sergipanos poderão escolher, entre diferentes trajetos, a formação \textbf{que} mais se harmoniza com suas pretensões e aptidões e com seu projeto de vida, a partir de \textbf{um} currículo \textbf{embasado} na formação integral.
Consideramos o espaço da escola \textbf{um} terreno fértil e \textbf{qualificado} para orientar os estudantes a se reconhecerem como \textbf{sujeitos}, considerando suas possibilidades e a relevância da participação e intervenção social na concretização de sua formação cidadã.
A ampliação da \textbf{percepção} das possibilidades para o futuro é fundamental para garantir o sucesso na construção do Projeto de Vida dos jovens.
Para tanto, é importante destacar o papel do desenvolvimento das habilidades socioemocionais atreladas ao Projeto de Vida no âmbito da Educação Básica, no desenrolar das práticas pedagógicas e, especificamente no Ensino Médio, tanto na Formação Geral Básica quanto nos Itinerários Formativos.
As Competências Gerais da BNCC aparecem no capítulo introdutório do documento e foram definidas a partir dos direitos éticos, estéticos e políticos assegurados pelas Diretrizes Curriculares Nacionais e dos conhecimentos, habilidades, atitudes e valores essenciais para a vida no \textbf{século} XXI. \textbf{Aqui} destacamos a competência 6 – Trabalho e Projeto de Vida, a qual se traduz em valorizar e apropriar -se de conhecimentos e experiências para entender o mundo do trabalho e fazer escolhas alinhadas à cidadania e ao seu projeto de vida com liberdade, autonomia, criticidade e responsabilidade.
Por sua vez, essa competência traz as subdimensões : Determinação, Esforço, Autoeficácia, Perseverança, Autoavaliação, Compreensão sobre o mundo do trabalho, Preparação para o trabalho.
Tomando esses elementos como substanciais para a formação integral dos jovens, as Ciências Humanas e Sociais Aplicadas lançam mão de \textbf{um} repertório próprio de objetos de conhecimento \textbf{que}, aliados ao uso de metodologias mais ativas e diversificadas, mobilizam as competências e habilidades específicas, advindas das competências gerais da BNCC, objetivando a promoção de situações de aprendizagem \textbf{baseadas} em experiências, contribuindo para a ampliação da capacidade de lidar com os conflitos pessoais e do contexto dos estudantes, no desenvolvimento do protagonismo e no seu amadurecimento para a tomada de decisões no \textbf{que} \textbf{diz} respeito aos mais diversos assuntos e trajetórias possíveis no curso e após a conclusão da educação básica.
Considerando esses pressupostos, e em articulação com as competências gerais da Educação Básica e com as da área de Ciências Humanas do Ensino Fundamental, as CHSA no Ensino Médio devem garantir aos estudantes o desenvolvimento de competências específicas e, relacionadas a cada \textbf{uma} delas, são indicadas, posteriormente, as habilidades a serem alcançadas nessa etapa.
O Currículo de Sergipe incorpora os Objetos de Conhecimento de cada Componente da Área ( Filosofia, Geografia, História e Sociologia ), em consonância com a BNCC, com o Novo Ensino Médio e todos os documentos correspondentes.
Importante destacar \textbf{que}, na nossa propositura, as Ciências Humanas e Sociais Aplicadas devem ser consideradas como \textbf{um} conjunto de objetivos de aprendizagem \textbf{que} se complementam para contribuir no projeto de vida e na formação integral dos jovens para a vida em sociedade e para a realização individual.
Nesse sentido, cumprem importante papel na garantia dos direitos de aprendizagem e desenvolvimento \textbf{que} fundamentam a nova composição do Ensino Médio, a saber : a Formação Geral Básica e os Itinerários Formativos.
Na Formação Geral Básica, deve ficar claro \textbf{que} Sergipe mantém os Componentes Curriculares \textbf{que} constituem a área, num movimento \textbf{que} pretende promover a \textbf{interdisciplinaridade}, a \textbf{contextualização} e a ocupação do espaço próprio das CHSA na formação integral dos jovens, em sintonia com as Diretrizes Curriculares Nacionais Gerais para a Educação Básica e para o Ensino Médio.
No intuito de atender aos princípios éticos, políticos e estéticos \textbf{que} fundamentam a BNCC e \textbf{que} devem estar presentes na Formação Geral Básica, apresentamos \textbf{um} Organizador Curricular pautado no \textbf{que} é essencial, dentro da carga horária destinada a cada Componente dentro da área, para a formação dos estudantes, contribuindo para o desenvolvimento da responsabilidade em relação aos direitos humanos, ao meio ambiente e à sua própria coletividade, tomando como parâmetro a justiça social.
Compete ainda às CHSA promover a formação de estudantes \textbf{que} articulem as \textbf{categorias} de pensamento histórico, geográfico, filosófico e sociológico em \textbf{face} de seu próprio tempo, por meio da observância e reflexão sobre as experiências humanas, em diferentes tempos, espaços e culturas e sob diversas lógicas de pensamento.
Com \textbf{face} nesse propósito, apresentamos também o Itinerário Formativo de Ciências Humanas e Sociais Aplicadas, destinado aos estudantes \textbf{que} escolherem se aprofundar nessa área, na perspectiva do Novo Ensino Médio, estruturado a partir de \textbf{um} percurso com começo, meio e fim, cujo fluxo \textbf{atravessa} os quatro eixos estruturantes e permite aos estudantes se desenvolverem de forma integral, orgânica e progressiva, lidando com desafios cada vez mais complexos.
Balizadas pela BNCC e pelos Referenciais Curriculares para elaboração de Itinerários Formativos, as 7 ( sete ) Atividades Integradoras foram elaboradas considerando também as seguintes questões propostas : O \textbf{que} os jovens \textbf{que} escolhem aprofundar e ampliar seus conhecimentos em CHSA precisam aprender nesse Itinerário Formativo, em articulação com a Formação Geral Básica?
Como podemos trazer tais conhecimentos por meio de objetos de conhecimento \textbf{que} entrelacem a Filosofia, Geografia, História e Sociologia, com vistas a \textbf{uma} \textbf{abordagem} interdisciplinar \textbf{que} estará sob a batuta de qualquer professor de \textbf{um} dos 4 ( quatro ) componentes, em qualquer unidade escolar de Sergipe?
A partir desses elementos \textbf{norteadores}, o Itinerário Formativo de CHSA traz a seguinte configuração : | ATIVIDADE INTEGRADORA | CH–MÓDULOS/AULA⁹ | |---|---:| | Núcleo de Estudos da África, da Diáspora e dos Povos Indígenas. | 80m/a | | Estudos Sergipanos. | 120m/a | | Observatório das Juventudes | 80m/a | | Laboratório de Estudos sobre Urbanização, Cidade e Meio Ambiente. | 40m/a | | Observatório de Mídias e Democracia¹⁰. | 100m/a | | Simulador da ONU. | 100m/a | | Cidadania e Direitos Humanos no Brasil. | 80m/a | ⁹ A carga horária apresentada para cada Atividade Integradora tem \textbf{um} caráter de sugestão pois, em virtude da autonomia dos professores para elaborar e desenvolver os Itinerários Formativos no âmbito da Unidade Escolar, essa configuração pode ser modificada, desde \textbf{que} sejam respeitados os preceitos do Novo Ensino Médio presentes em toda legislação vigente, nos documentos oficiais, Matriz Curricular e no Currículo de Sergipe – Etapa Ensino Médio. ¹⁰ Essa Atividade Integradora, além de compor o Itinerário de Aprofundamento em CHSA, tem como possibilidade ser \textbf{um} repositório de toda a produção dos estudantes da escola, especialmente dos Itinerários Formativos de todas as áreas do conhecimento.
Pois bem, nossa elaboração percorreu \textbf{um} caminho pautado na concepção da \textbf{interdisciplinaridade} a partir \textbf{um} currículo elaborado por integração.
Isso quer \textbf{dizer} \textbf{que} não há delimitações estanques \textbf{que} se somam para formar \textbf{um} mosaico de objetos próprios de cada componente.
Nossa proposta traz objetos de conhecimento relacionados e de forma mais aberta, dilui as \textbf{fronteiras} disciplinares rígidas, sua organização propicia experiências de aprendizagem mais formativas através de metodologias ativas e permite \textbf{uma} visão mais ordenada e integrada do mundo.
Trazemos temas essenciais a serem trabalhados pelo professor de CHSA \textbf{que} estará mediando as Atividades Integradoras, nos formatos \textbf{que} foram escolhidos como mais apropriados para mobilizar as habilidades ligadas ao(s) Eixo(s) Estruturante(s) correspondente(s ).
Os planos de atividade docente são propostas estruturadas para a elaboração dos planos de curso \textbf{que} abrem espaço para a autonomia e a criatividade pedagógica de cada professor da Rede Pública Estadual de Sergipe.
Entendemos \textbf{que}, antes de mobilizar habilidades ligadas ao fazer criativo, à atuação sociocultural e ao empreendedorismo em Ciências Humanas e Sociais Aplicadas, é \textbf{imprescindível} aguçar, aprofundar o fazer científico, investigativo, à luz do repertório produzido por essa área.
Sendo assim, organizamos o Itinerário Formativo de modo \textbf{que}, em todas as Atividades Integradoras, o estudante começa se debruçando sobre a Investigação Científica e, num segundo momento, ele faz as devidas associações com os demais Eixos Estruturantes, balizados pela lente das Ciências Humanas e Sociais Aplicadas.
Por fim, outro ponto relevante \textbf{que} foi levado em consideração nessa construção foi a associação com os temas \textbf{contemporâneos} transversais da BNCC, a saber : 1.
Meio Ambiente – Educação Ambiental e Educação para o consumo ; 2.
Economia – Trabalho, Educação Financeira e Educação Fiscal ; 3.
Saúde – Saúde, Educação Alimentar e Nutricional ; 3.
Cidadania e Civismo – Vida familiar e social, Educação para o trânsito, Educação em Direitos Humanos, Direitos da Criança e do Adolescente, Processo de envelhecimento, respeito e valorização do idoso ; 4.
Multiculturalismo – Diversidade Cultural, Educação para a valorização do multiculturalismo nas matrizes históricas e culturais brasileiras ; 5.
Ciência e Tecnologia.
Eles perpassam as Atividades Integradoras, dialogando com as habilidades, unidades curriculares e temas propostos.
A partir dessa proposta, esperamos \textbf{que} as Ciências Humanas e Sociais Aplicadas, em articulação com as outras áreas do conhecimento e outros saberes \textbf{que} extrapolam a \textbf{excelência} acadêmica proporcionem aos estudantes a capacidade de realizar pesquisas científicas, criar obras, soluções e/ou inovações, intervir positivamente na realidade e empreender iniciativas pessoais, acadêmicas, produtivas e/ou cidadãs, sempre em diálogo com as Competências Gerais indicadas pela BNCC como finalidade da Educação Básica e direito de aprendizagem e desenvolvimento de todos os jovens.
A seguir, elucidamos as competências específicas para CHSA no Ensino Médio e, em dois quadros demonstrativos, o percurso de competências e habilidades gerais e da área, bem como dos eixos estruturantes, para a elaboração do currículo de CHSA – Formação Geral Básica e Itinerário Formativo, e apresentamos cada componente da área. 1.
Analisar processos políticos, econômicos, sociais, ambientais e culturais nos âmbitos local, regional, nacional e mundial em diferentes tempos, a partir da pluralidade de procedimentos epistemológicos, científicos e tecnológicos, de modo a compreender e posicionar -se criticamente em relação a eles, considerando diferentes pontos de vista e tomando decisões \textbf{baseadas} em argumentos e fontes de natureza científica. 2.
Analisar a formação de territórios e \textbf{fronteiras} em diferentes tempos e espaços, mediante a compreensão das relações de poder \textbf{que} determinam as \textbf{territorialidades} e o papel geopolítico dos Estados-nações. 3.
Analisar e avaliar criticamente as relações de diferentes grupos, povos e sociedade com a natureza ( produção, distribuição e consumo ) e seus impactos econômicos e socioambientais, com vistas à proposição de alternativas \textbf{que} respeitem e promovam a consciência, a ética socioambiental e o consumo responsável em âmbito local, regional, nacional e global. 4.
Analisar as relações de produção, capital e trabalho em diferentes territórios, contextos e culturas, discutindo o papel dessas relações na construção, consolidação e transformação das sociedades. 5.
Identificar e combater as diversas formas de injustiças, preconceito e violência, adotando princípios éticos, democráticos, inclusivos e solidários, e respeitando os Direitos Humanos. 6.
Participar do debate público de forma crítica, respeitando diferentes posições e fazendo escolhas alinhadas ao exercício da cidadania e ao seu projeto de vida, com liberdade, autonomia, consciência crítica e responsabilidade.

% matched lemmas: abordagem, aqui, atravessar, atual, basear, categoria, compromisso, conceitual, contemporâneo, contextualização, disciplina, dizer, embasar, enquanto, erro, excelência, face, focar, fronteira, globalizado, historicamente, humanidade, implicar, imprescindível, interdisciplinaridade, norteador, percepção, problematização, qualificar, que, raciocínio, sistematização, sujeito, século, territorialidade, teórico, um
\end{textsample}
